\documentclass[a4paper,11pt]{article}
\usepackage[italian]{babel}
\usepackage{geometry}
\geometry{margin=2.5cm}

\title{Soluzioni – Verifica SQL ITS}
\date{}

\begin{document}
	\maketitle
	
	\begin{enumerate}
		
		\item SELECT * FROM STUDENTE;
		
		\item SELECT nome,cognome
		FROM STUDENTE
		WHERE annoIscrizione>2022;
		
		\item SELECT *
		FROM AZIENDA
		WHERE settore='Software';
		
		\item SELECT *
		FROM PROGETTO STAGE
		WHERE durataMesi>6;
		
		\item SELECT S.nome,S.cognome,A.ragioneSociale
		FROM STUDENTE S
		JOIN STAGE ST ON S.idStudente=ST.idStudente
		JOIN PROGETTO STAGE P ON ST.idProgetto=P.idProgetto
		JOIN AZIENDA A ON P.idAzienda=A.idAzienda;
		
		\item SELECT DISTINCT idStudente
		FROM STAGE;
		
		\item SELECT *
		FROM STUDENTE
		WHERE idStudente NOT IN
		(SELECT idStudente FROM STAGE);
		
		\item SELECT A.ragioneSociale,COUNT(*)
		FROM AZIENDA A
		JOIN PROGETTOSTAGE P ON A.idAzienda=P.idAzienda
		JOIN STAGE S ON P.idProgetto=S.idProgetto
		GROUP BY A.ragioneSociale;
		
		\item SELECT idStudente,AVG(valutazione)
		FROM STAGE
		GROUP BY idStudente;
		
		\item SELECT A.ragioneSociale
		FROM AZIENDA A
		JOIN PROGETTOSTAGE P ON A.idAzienda=P.idAzienda
		JOIN STAGE S ON P.idProgetto=S.idProgetto
		GROUP BY A.ragioneSociale
		HAVING COUNT(*)>3;
		
		\item SELECT idTutorAz,COUNT(*)
		FROM STAGE
		GROUP BY idTutorAz;
		
		\item SELECT C.nomeCompetenza
		FROM STUDENTECOMPETENZA SC
		JOIN COMPETENZA C
		ON SC.idCompetenza=C.idCompetenza
		WHERE idStudente=1;
		
		\item SELECT P.titolo
		FROM PROGETTOSTAGE P
		JOIN PROGETTOCOMPETENZA PC
		ON P.idProgetto=PC.idProgetto
		JOIN COMPETENZA C
		ON PC.idCompetenza=C.idCompetenza
		WHERE C.nomeCompetenza='Java';
		
		\item SELECT DISTINCT S.nome,S.cognome
		FROM STUDENTE S
		JOIN STUDENTECOMPETENZA SC
		ON S.idStudente=SC.idStudente
		JOIN COMPETENZA C
		ON SC.idCompetenza=C.idCompetenza
		WHERE C.nomeCompetenza='SQL';
		
		\item SELECT idStudente
		FROM STAGE ST
		JOIN PROGETTOSTAGE P
		ON ST.idProgetto=P.idProgetto
		JOIN AZIENDA A
		ON P.idAzienda=A.idAzienda
		GROUP BY idStudente
		HAVING COUNT(DISTINCT A.città)>1;
		
		\item SELECT *
		FROM AZIENDA
		WHERE idAzienda NOT IN
		(SELECT idAzienda FROM PROGETTOSTAGE);
		
		\item SELECT idAzienda,AVG(durataMesi)
		FROM PROGETTOSTAGE
		GROUP BY idAzienda;
		
		\item SELECT idStudente
		FROM STAGE
		GROUP BY idStudente
		HAVING AVG(valutazione)>27;
		
		\item SELECT idTutorAcc
		FROM STAGE
		GROUP BY idTutorAcc
		ORDER BY COUNT(*) DESC
		LIMIT 1;
		
		\item SELECT idProgetto,COUNT(idStudente)
		FROM STAGE
		GROUP BY idProgetto;
		
		\item SELECT idCompetenza
		FROM PROGETTOCOMPETENZA
		GROUP BY idCompetenza
		ORDER BY COUNT(*) DESC
		LIMIT 1;
		
		\item SELECT idStudente
		FROM STUDENTECOMPETENZA
		GROUP BY idStudente
		HAVING COUNT(idCompetenza)=
		(SELECT COUNT(*) FROM COMPETENZA);
		
		\item SELECT idStudente
		FROM STAGE
		GROUP BY idStudente
		HAVING COUNT(DISTINCT idTutorAz)>1;
		
		\item SELECT MAX(valutazione)
		FROM STAGE;
		
		\item SELECT idStudente
		FROM STAGE
		GROUP BY idStudente
		ORDER BY AVG(valutazione) DESC
		LIMIT 1;
		
	\end{enumerate}
	
	CREATE TABLE AULA(
	idAula INT PRIMARY KEY,
	edificio VARCHAR(50) NOT NULL,
	capienza INT NOT NULL CHECK(capienza>0)
	);
	
	CREATE VIEW STUDENTIVALUTAZIONE AS
	SELECT S.idStudente,S.nome,S.cognome,
	AVG(ST.valutazione) AS mediaValutazioneStage
	FROM STUDENTE S
	JOIN STAGE ST
	ON S.idStudente=ST.idStudente
	GROUP BY S.idStudente,S.nome,S.cognome;
	
\end{document}
